\documentclass[11pt,a4paper]{article}
\usepackage{fshf-style}
\SetFshfShortTitle{Fixed Income Portfolio}

\begin{document}

\makefshftitle{Fixed Income Portfolio}%
{Fixed Income --- EUR and USD Sovereign, Corporate and Floating Rate}%
{Anastasia Shevchuk, Johannes Volkemer \& Raphael Banner}%
{2 July 2026}%
{\relax}

\setcounter{tocdepth}{2}
\tableofcontents
\clearpage

\section{Goals and Scope of work}

\subsection{Introduction}

This report examines a strategy for allocating to fixed-income exchange-traded funds across the Euro and USD currency areas, against the backdrop of current and anticipated monetary policy decisions by the European Central Bank and the US Federal Reserve. The central question is how a conservatively positioned fixed income portfolio can be structured to respond to the changed interest rate environment without incurring excessive risk.

The analysis covers information available up to and including 26 June 2026. Political and geopolitical developments after this cut-off date, including the renewed escalation of tensions related to the Iran conflict, are deliberately excluded. This methodological boundary is noted again as a limitation in the concluding section.

Methodologically, the report combines a macroeconomic analysis of monetary policy conditions in both currency areas in Section 2 with a derivation of strategic implications in Section 3 and a structured product selection in Section 4. Products are evaluated according to criteria including interest rate sensitivity, credit quality, liquidity and replication method.

\subsection{Disclaimer}

This report was produced as part of an academic project and is intended solely for academic purposes within the context of education and research. All statements, assessments and product selections do not constitute investment advice and must not be construed as a recommendation or solicitation to buy or sell any financial instrument. The authors are not licensed investment advisers.

All market data, interest rates and product metrics reflect the state of knowledge at the time indicated and may have changed since. No warranty is given as to the completeness or accuracy of the information presented. Before making any actual investment decision, readers are advised to consult current product documentation, including the key investor information document and prospectus of the relevant funds, as well as independent professional advice.

\subsection{Our Mandate}

The mandate is built around three guiding principles. First, the portfolio is designed to provide a safe foundation and to act as a shock absorber within a broader investment context, dampening volatility rather than generating additional risk. Second, the portfolio deliberately targets a lower sensitivity to interest rate movements relative to long-duration fixed income benchmarks, referred to here as a lower beta profile, in order to reduce dependence on any single monetary policy driver. Third, this reduction in single-factor dependency is achieved through deliberate diversification, both across the two currency areas of Euro and USD and across different maturity segments through a ladder of duration buckets.

Consistent with the portfolio construction framework defined in Section 3.3, the EUR and USD sleeves are weighted at 70\% and 30\% of total portfolio value respectively. The EUR sleeve follows the mild barbell structure with a moderate front-end bias established in Section 3.1, comprising 50\% in short duration, 15\% in intermediate duration and 35\% in long duration, while the USD sleeve follows the bullet-weighted, intermediate-focused structure established in Section 3.2, comprising 20\% in short duration, 50\% in intermediate duration and 30\% in long duration.

The mandate is active rather than purely passive. Against the backdrop of a volatile political and monetary policy environment, FS Student Hedge Fund retains the ability to adjust weights between maturity segments and between the two currency areas over time through rebalancing. Specific triggers for such adjustments are identified in Section 5.

\section{Analysis}

\subsection{Global Context}

As noted in Section 1.1, the analysis covers developments up to and including 26 June 2026. Further geopolitical developments are not incorporated into our analysis.

\subsection{Eurozone}

\subsubsection{Growth and Inflation Backdrop}

The Eurozone is currently best described as a low-growth environment with clear inflation risks. Growth momentum is on the weaker side, with euro area GDP contracting by \(-\)0.2\% q/q in Q1 2026, while inflation re-accelerated to 3.2\% y/y in May 2026\footnote{European Commission, Eurostat (2026). GDP down by 0.2\% and employment up by 0.1\% in the euro area. Euro Indicators. Available at: \url{https://ec.europa.eu/eurostat/web/products-euro-indicators/w/2-05062026-ap}}. Against this backdrop, the ECB raised all three key policy rates by 25bp on 11 June 2026, taking the deposit facility rate to 2.25\%, the MRO rate to 2.40\% and the marginal lending facility rate to 2.65\%\footnote{European Central Bank (2026). Monetary Policy Decisions, 11 June 2026. Available at: \url{https://www.ecb.europa.eu/press/pr/date/2026/html/ecb.mp260611~4d41bd5e83.en.html}}\textsuperscript{,}\footnote{National Bank of Belgium (2026). Monetary Policy Decisions, June 2026. Available at: \url{https://www.nbb.be/en/news-events/news/monetary-policy-decisions-june-2026}}. The main tension is that the ECB is tightening into weak growth because upside inflation risks remain material. Higher energy prices could lift headline inflation and, if sustained, feed into wages, services and core inflation. Market pricing currently implies roughly one further ECB hike this year, with one pricing source showing the deposit rate around 2.50\% by December 2026. This should be treated as market-implied pricing rather than a firm forecast or an ECB signal\footnote{Rate Probability (n.d.). ECB Rate Probability Tracker. Available at: \url{https://rateprobability.com/ecb}}\textsuperscript{,}\footnote{ECB Watch (n.d.). ECB Interest Rate Expectations and Monetary Policy Monitoring. Available at: \url{https://ecb-watch.eu/}}. The core thesis is that Eurozone rates are no longer trading a simple disinflation or easing story; markets are pricing renewed inflation risk despite a weak growth backdrop.

\subsubsection{Geopolitical Layer}

The key geopolitical issue for Eurozone fixed income is the Middle East energy shock, because it affects inflation through oil, gas, electricity and broader production costs. The ECB cited the war in the Middle East as a source of inflation pressure when it raised rates in June 2026, although this should be treated as one important factor rather than the sole reason for the decision. The main transmission channels are higher headline inflation, higher production costs, possible second-round effects in core goods and services, weaker real incomes, lower confidence and higher risk premia. Russia's war on Ukraine remains highly relevant for energy security, defense spending, fiscal pressure and risk sentiment, but the immediate market focus has shifted toward the Middle East because of its current impact on energy prices and inflation assumptions. For fixed income, geopolitical risk should be assessed through Bund flight-to-quality potential, inflation breakevens, short-end ECB repricing, peripheral spread risk and higher term premia linked to fiscal and defense spending.

\subsubsection{Macroeconomic indicators}

At the same time, inflation remains above target and continues to constrain monetary policy. Euro area inflation increased to 3.2\% y/y in May 2026, up from 3.0\% in April\footnote{European Commission, Eurostat (n.d.). Euro Indicators. Available at: \url{https://ec.europa.eu/eurostat/news/euro-indicators}}. The ECB's June projections put headline inflation at 3.0\% in 2026, 2.3\% in 2027 and 2.0\% in 2028, while inflation excluding energy and food (i.e.\ core inflation) is projected at 2.5\% in both 2026 and 2027 and 2.2\% in 2028. Core inflation is therefore the more important signal for the ECB than headline inflation alone. The key risk is that an energy shock does not remain isolated, but gradually feeds into production costs, wages and services prices. This is why the ECB is unlikely to treat the current inflation overshoot as purely temporary unless incoming data shows a clear decline in underlying inflation momentum.

\fshffig{fi-fig1-annual-inflation-rates.png}{Annual inflation rates}{Source: European Central Bank (n.d.). ECB Data Portal. Available at: \url{https://data.ecb.europa.eu/}}

The labor market reinforces this policy dilemma. The ECB Data Portal shows 6.3\% unemployment rate\footnote{European Central Bank (n.d.). ECB Data Portal. Available at: \url{https://data.ecb.europa.eu/}}. Labor cost pressure also remains relevant, with hourly labor costs rising by 3.2\% y/y in Q1 2026 in the euro area. The ECB continues to monitor negotiated wages closely because wage growth is a key driver of services inflation persistence\footnote{European Central Bank (n.d.). ECB Wage Tracker (EWT) Dataset. Available at: \url{https://data.ecb.europa.eu/data/datasets/EWT}}. Taken together, the macro data do not point to a clean disinflationary environment. Growth is weak, but inflation and wages remain firm enough to keep the ECB cautious. For a fixed income portfolio, this argues against relying solely on a growth-led rally in bonds and supports a more balanced duration approach.

\subsubsection{ECB Framework}

The ECB's reaction function is data-dependent and meeting-by-meeting, with no pre-commitment to a specific rate path. Policy decisions are based on the inflation outlook, risks around inflation, incoming economic and financial data, underlying inflation dynamics and the strength of monetary policy transmission. In the current environment, this means the ECB has to balance weak growth against the risk that the energy shock becomes embedded in wages and services inflation. The June hike suggests that inflation persistence and second-round effects currently carry more weight than near-term growth weakness. This is also supported by the ECB's own communication, which framed the 25bp hike as a response to inflation pressures from the Middle East energy shock, while stressing that the outlook remains uncertain and that the Governing Council is not pre-committing to a specific path. If inflation stays sticky, the front end could remain vulnerable to further hawkish repricing. If growth deteriorates sharply, the curve could bull-steepen as markets price future easing again.

\subsection{USA}

\subsubsection{US Research Methodology}

The Federal Reserve Bank (Fed) is the central bank for the United States of America (USA) and primarily responsible for ensuring financial stability with its monetary policies\footnote{Federal Reserve Board (n.d.). Monetary Policy. Available at: \url{https://www.federalreserve.gov/aboutthefed/fedexplained/monetary-policy.htm}}. Through adjustments of the federal funds rate and its balance sheet, the Fed influences short-term interest rates, financial conditions for the economy as well as dynamics of inflation by reacting to and continuously assessing economic indicators such as the consumer price index (CPI)\footnote{Federal Reserve Board (n.d.). PCE Inflation. Available at: \url{https://www.federalreserve.gov/economy-at-a-glance-inflation-pce.htm}}. Changes in the monetary policy of the Fed translate to impacts across the entire Treasury Yield curve and ultimately affect the valuation of fixed-income assets. Consequently, understanding the Fed's reaction and functionalities is essential for our report to assess economic forecasts of future interest rate developments.

In doing so, the Fed not only focuses on pure inflation-related data like the European Central Bank (ECB) but also takes employment data into account. This dualistic mandate clearly separates it from the ECB and has significant implications for its interest rate decisions.

The first mandate is to ensure price stability and preserve the buying power of the US Dollar. In practice, the long-term target range for price inflation lies approximately around 2\%, to avoid high-inflation scenarios and simultaneously stimulate the economy and consumer spending. Periods of deflation are undesirable, as these may weaken demand and construct a downward spiral of increasing downward pressure on prices to catch up with consumer demand.

The second mandate of the Fed is formulated as to promote the highest possible level of employment while also preserving price stability. Although there is no defined target range for employment, the Fed tries to prevent the labor market from overheating and to ensure a ``soft landing'' in case of economic downturns. Labor market conditions are assessed through a broad range of economic indicators, including but not limited to unemployment rates, payroll growth, new job openings, or overall wage growth.

As those two mandates are served in parallel and have the same priority to the Fed, this implies that both economic dimensions must be taken into account when assessing the likelihood for interest rate developments. We therefore widen our research approach to include labor market data for our US research part.

The Fed mainly relies on economic data to decide on their monetary policy reactions. A typical classification of economic indicators is done in three categories:
\begin{enumerate}[label=\Roman*.]
\item Leading indicators
\item Coincident indicators
\item Lagging indicators
\end{enumerate}

Policymakers evaluate a broad set of these economic indicators to assess the past, present, and expected future developments of the economy and related conditions\footnote{Federal Reserve Bank of St.\ Louis (n.d.). FRED --- Federal Reserve Economic Data. Available at: \url{https://fred.stlouisfed.org}}. We therefore also assess these three types of indicators, to develop a comprehensive view and expectation for our recommendations.

\subsubsection{Analysis of Economic Indicators}

\paragraph{Inflation and Price Stability Assessment}

The first economic assessment of the US is done looking at the first mentioned mandate of the Fed, being moderate inflation and price stability. In this section, we assess the past and current state of price levels to assess potential implications for the interest rate set by the Fed.

\fshffig{fi-fig2-core-inflationary-data.png}{Core inflationary data}{Source: FRED, series PCEPILFE, percent change from year ago}

As seen in the graph, core inflationary indicators such as the Personal Consumption Expenditures (PCE; Blue) and the Sticky Price Consumer Price Index (CPI; Green) are seen to be correlated and tracking each other's movements, with the Sticky Price CPI lagging roughly one year behind. Both indicators exclude energy and food prices, as those are seen to be very volatile, especially with recent geopolitical and macroeconomic developments around the world.

With the peak in core PCE of approximately 5.6\% happening in February 2022, the level of core inflation declined to approximately 3.1\% by the end of 2023 (revised data; the initially reported reading was closer to 2.9\%), remaining above the mentioned target inflation level for the Fed, and bottomed near 2.6\% in early 2025. From there, inflation data shows a renewed increase up until the most recent level of approximately 3.3\% y/y in April 2026. In the same period, the Fed Funds rate target range declined from a high of 5.25\%--5.50\%, set in July 2023 and held into August 2024, to 3.50\%--3.75\% in May 2026\footnote{Federal Reserve Bank of New York (n.d.). Effective Federal Funds Rate (EFFR). Available at: \url{https://www.newyorkfed.org/markets/reference-rates/effr}}\textsuperscript{,}\footnote{CME Group (n.d.). CME FedWatch Tool. Available at: \url{https://www.cmegroup.com/markets/interest-rates/cme-fedwatch-tool.html}}.

\paragraph{Labor Market Assessment}

In the second dimension of the Fed mandate, we assess the past and current labor market conditions to gain a full understanding of the economic state of the US. For the Fed, especially indicators towards employment and wage growth contain helpful information to gauge the current sentiment and economic outlooks.

\fshffig{fi-fig3-core-labor-market-data.png}{Core labor market data}{Source: FRED. Available at: \url{https://fred.stlouisfed.org/graph/?g=1WWDy}}

The unemployment rate (UR; Green) is the primary real-time gauge for the Fed to assess the current state of the labor market. Following the dualistic mandate, the Fed tries to maintain the highest level of employment while preserving price stability. Turning points in the UR are consistently preceding or coinciding with economic downturns, as observable in the graph, especially drastic in 2020. After recent shocks from the Covid-crisis, which exhibited sharp increases in unemployment, the labor market stabilized from 2022 onwards, coming back to pre-covid levels. A modest rise since 2023 suggests a gradual normalization to desired levels rather than a sharp implication of downturns in the labor market.

As seen in the graph, the overall Fed Funds Rate (Red; dotted) is also linked to both Sticky Price CPI and WG, lagging around one year behind. This shows the dualistic nature of the Fed and its mandate, both driven by inflation levels and the labor market. Sharp increases in the UR are responded with a cut in interest rates, while recent decreasing tendencies are responded with modest decreases in the Funds target range. On the other side, interest rate hikes can be observed after periods of higher WG, in connection with higher inflation levels, typically lagging one year behind.

\paragraph{Market Expectations and Implications for Federal Reserve Interest Rate Policy}

As seen on both assessments, inflationary tendencies tend to persist and might pick up again in the near future. Whether it stems from personal consumption or wage growth, there is imminent inflationary pressure. Both assessments show, assessing the full state of the economy is not possible with only looking at single indicators but rather taking the interplay between them into consideration.

To further back up our assessments, we want to take policymakers' personal perspective of potential interest rate developments into account.

\fshffig{fi-fig4-fomc-dot-plot.png}{``Dot Plot'' of FOMC members --- Summary of Economic Projections, June 17, 2026}{Note: Each shaded circle indicates the value (rounded to the nearest 1/8 percentage point) of an individual participant's judgment of the midpoint of the appropriate target range for the federal funds rate or the appropriate target level for the federal funds rate at the end of the specified calendar year or over the longer run.}

The figure shows the famous ``Dot Plot'' of the FED and its individual policymakers, which is widely used to assess probabilities of further hikes or cuts in interest rates. The figure shows the FED Funds Rate expectation of each member for the next full three years until 2028 and a longer run projection. In red, the current level of the FED Funds Rate is highlighted, showing a majority of members having their estimate for the year 2026 at the current range, with a slight tilt upwards. Following the years, members overall expect decreasing levels of FED Funds Rates, with slight uncertainties and mixed views for the next years. These projections are in concordance with our assessments of the price levels and the labor market, as those two indicate at least stable FED Funds Rates for the ongoing year and rate cuts to happen as early as 2027 onwards. This is in line with a major research report recently published by Goldman Sachs, which revised its interest rate cut expectations from end of year 2026 to mid- and end-year 2027, as long as there are no further supply shocks hitting pricing levels.

\section{Implications and construction of our portfolio}

\subsection{Eurozone}

\subsubsection{Rates Market Implications}

With one further ECB hike already priced, the duration view depends on whether inflation persistence is adequately reflected in current yields. If inflation remains sticky or the ECB needs to hike more than expected, duration exposure should stay cautious, particularly at the front end where yields are most directly linked to policy-rate expectations. If growth weakness becomes the dominant driver, longer duration becomes more attractive because yields could decline and high-quality bonds could regain their defensive role. The front end already embeds some hawkishness, but it is not fully protected against further upside surprises in inflation or wages\footnote{ECB Watch (n.d.). ECB Interest Rate Expectations and Monetary Policy Monitoring. Available at: \url{https://ecb-watch.eu/}}. The long end is driven less by the next ECB meeting and more by term premia, fiscal supply, energy risk and global rates. The ECB Data Portal shows the euro area AAA-rated 10-year government bond yield at 2.92\% on 26 June 2026 and is visualized below in the chart\footnote{European Central Bank (n.d.). ECB Data Portal. Available at: \url{https://data.ecb.europa.eu}}. The curve outlook therefore depends on the type of shock. Short-end ECB repricing would point to bear flattening, term-premium or fiscal concerns would point to bear steepening, and a growth shock would support bull steepening. Breakevens should remain sensitive to energy-linked inflation volatility, especially if markets become concerned about second-round effects in wages and services. Since the front end already reflects one additional ECB hike, significant repricing risk appears limited under the base case, but the portfolio should still avoid extreme positioning in short maturities.

\fshffig{fi-fig5-aaa-10y-yield.png}{Euro area AAA-rated 10-year government bond yield}{Source: European Central Bank (n.d.). ECB Data Portal. Available at: \url{https://data.ecb.europa.eu}}

\subsubsection{Sovereign Spreads}

Eurozone sovereign spreads should be analyzed through growth divergence, fiscal sustainability, ECB backstops and market risk appetite. Fiscal risk remains important because the euro area government debt ratio stood at 87.8\% of GDP in Q4 2025. Country dispersion is material: Germany's 2025 debt ratio was 63.5\% of GDP, France 115.6\%, Italy 137.1\%, Spain 100.7\% and Greece 146.1\%, based on Eurostat data shown by Destatis\footnote{Federal Statistical Office of Germany (Destatis) (n.d.). Economy and Finance --- Key Indicators for Europe. Available at: \url{https://www.destatis.de/Europa/EN/Topic/Key-indicators/EconomyFinance.html}}. Italy and France require close monitoring because high debt ratios and higher rates increase sensitivity to refinancing costs, rating risk and political risk. Spread pressure also depends on maturity profiles, primary balances, ECB backstops, investor base and domestic politics. ECB tools such as PEPP flexibility and the Transmission Protection Instrument remain relevant, but they do not remove spread risk if fiscal credibility weakens. Given the portfolio's focus on stability and low volatility, spread exposure should remain controlled.

\subsubsection{Portfolio Construction \& Positioning}

Against this backdrop, a balanced duration profile with a moderate front-end bias and a mild barbell structure appears appropriate with an allocation of 50\% short duration, 15\% intermediate duration and 35\% long duration. The short-duration allocation limits overall volatility, reduces sensitivity to unexpected ECB repricing and provides stable carry while preserving capital. The intermediate allocation keeps exposure across the curve and avoids making the portfolio too dependent on one specific curve outcome. The long-duration allocation preserves diversification against equity drawdowns and provides protection if growth deteriorates enough for yields to fall.

\subsubsection{Investment Conclusion}

The investment conclusion is that the Eurozone fixed income environment supports a defensive, balanced allocation rather than a strong directional duration position. Growth momentum has weakened, but inflation remains above the ECB's target and wage-related services inflation is still relevant. This keeps the ECB cautious and limits the case for positioning the portfolio purely for a growth-led bond rally. At the same time, because markets already price approximately one additional ECB hike, the central scenario does not require an aggressively short-duration stance either.

The preferred positioning is therefore a moderate front-end bias combined with selective long-duration exposure. The front-end allocation supports stability and limits overall volatility, while the long-duration allocation preserves diversification in adverse macroeconomic scenarios. This is important because fixed income should not only generate carry but also provide resilience if economic weakness becomes the dominant market driver. The portfolio should therefore avoid excessive exposure to both short-end ECB repricing risk and long-end duration risk.

Overall, a balanced fixed income strategy with a moderate front-end bias and a mild barbell structure appears most appropriate. Such an allocation prioritizes capital preservation, limits portfolio volatility and maintains diversification benefits. It also fits the broader objective of using fixed income as a stable foundation within a multi-asset portfolio rather than as a vehicle for taking large macro bets.

\subsection{United States}

\subsubsection{Interest rate market implications}

The US fixed-income market is characterized by a fundamentally different macroeconomic backdrop than the Eurozone. While growth remains more resilient, inflation continues to exceed the Federal Reserve's long-run target and labor market conditions remain sufficiently strong to prevent an imminent shift toward accommodative monetary policy. The analysis of both inflation and labor market indicators suggests that inflationary pressures have become more persistent than previously expected and that the Federal Reserve has limited room to further reduce interest rates in the near term.

The inflation assessment shows that both PCE inflation and Sticky Price CPI remain elevated, with historical patterns indicating that inflation persistence can last significantly longer than initial shocks suggest. Rising energy prices and geopolitical uncertainty further increase the risk that inflation remains above target throughout 2026. At the same time, labor market indicators do not signal a meaningful deterioration in economic conditions. Unemployment remains low by historical standards, while wage growth continues to support service-sector inflation. Together, these dynamics are consistent with a Federal Reserve that is likely to maintain restrictive monetary policy for an extended period.

The Federal Reserve's June 2026 Summary of Economic Projections reinforces this view. Most policymakers expect interest rates to remain largely unchanged during 2026, while the majority of projected rate cuts are concentrated in 2027 and beyond.

For duration positioning, this environment supports a balanced approach with a moderate overweight to intermediate maturities. Short-duration bonds continue to benefit from attractive yields and lower volatility while providing protection against further inflation-related repricing risk. Intermediate maturities offer an attractive balance between carry and duration exposure, particularly if the market gradually begins to price the first stages of a future easing cycle. Long-duration bonds on the other hand remain valuable as a portfolio diversifier and safety anchor against short-term volatility.

\subsubsection{Portfolio Construction}

Given the current macroeconomic environment and in comparison to the Eurozone results, the preferred USD fixed-income allocation consists of 20\% short duration, 50\% intermediate duration and 30\% long duration. This way, we can make use of later expected rate cuts in the US opposed to the Eurozone while at the same time maintain an anchor in our portfolio with longer-term exposure.

The 20\% short-duration allocation provides stability, preserves liquidity and limits sensitivity to potential upward repricing of policy expectations should inflation remain stubbornly high. Although the Federal Reserve is expected to keep policy restrictive, the likelihood of substantial additional rate hikes appears limited, reducing the need for an excessive short-duration bias.

Intermediate maturities represent the portfolio's core exposure. This segment offers an attractive combination of yield, moderate duration risk and participation in a potential future easing cycle. Given the expectation that policy rates remain stable throughout 2026 before gradually declining in 2027, intermediate maturities provide favorable risk-adjusted exposure to the most likely policy path.

The allocation towards longer-term durations serves primarily as a strategic hedge against growth disappointments, recession risks and broader risk-off market environments. While inflation persistence currently limits the upside potential of long-duration bonds, they remain an important source of diversification and portfolio protection should economic momentum weaken more significantly than anticipated.

\subsection{Currency Allocation (EUR vs.\ USD)}

Given the portfolio's objective of capital preservation and low volatility, the fixed income allocation should maintain a clear bias towards euro-denominated assets. While US dollar fixed income can provide diversification benefits through exposure to a different economic and monetary policy cycle, these benefits must be weighed against the additional risks arising from foreign exchange exposure.

For a euro-based investor, EUR-denominated bonds avoid currency mismatches between assets and liabilities and eliminate exchange-rate fluctuations as a source of portfolio volatility. Although USD bonds may offer attractive yields and access to the highly liquid US Treasury market, unhedged USD exposure can introduce return variability that is unrelated to the underlying credit quality or duration profile of the bond. In some periods, currency movements can be larger than the bond returns themselves.

Currency hedging can reduce this exchange-rate risk, but it comes at a cost. Hedging expenses are driven by interest rate differentials between the United States and the Eurozone and can materially reduce the yield advantage that initially makes USD bonds attractive. Hedging also adds operational complexity and increases implementation costs.

Given the portfolio's defensive mandate, the euro allocation should therefore remain the dominant component of the fixed income portfolio. A strategic allocation of approximately 70\% EUR fixed income and 30\% USD fixed income provides a reasonable balance between stability and diversification. The euro allocation serves as the portfolio's core anchor, while the smaller USD allocation provides exposure to a different interest-rate cycle without allowing exchange-rate risk to become a dominant driver of portfolio performance.

\section{Products}

The product selection in the EUR sleeve reflects the barbell structure defined in Section 3.1 (50\% short duration, 15\% intermediate duration, 35\% long duration), while the USD sleeve reflects the bullet-weighted intermediate-focused structure defined in Section 3.2 (20\% short duration, 50\% intermediate duration, 30\% long duration). The overall portfolio is split 70\% EUR and 30\% USD as established in Section 3.3. All products are UCITS-compliant and, unless otherwise noted, have been verified as tradeable through Interactive Brokers.

Within each bucket, weights follow each product's role as described in the function column. The position identified as the bucket's primary anchor receives the largest share. Laddering or transitional positions receive progressively smaller shares. Positions added specifically for provider or index diversification are sized as smaller satellite allocations rather than core building blocks.

\subsection{EUR Products (70\% of total portfolio)}

The EUR sleeve spans all three duration buckets of the barbell across 12 products. The short-duration allocation of 50\% is composed of money market instruments and floating rate bonds that absorb ECB rate moves without meaningful price impact, complemented by short-dated government bonds spread across multiple index providers to avoid benchmark concentration. The intermediate allocation of 15\% introduces spread income through investment grade corporate bonds and covered bonds. The long-duration allocation of 35\% provides diversification against equity drawdowns and preserves the portfolio's defensive character in scenarios of sharp economic deterioration, where high-quality long-dated government bonds historically benefit from flight-to-quality dynamics. This allocation also ensures the portfolio retains convexity should the ECB pivot earlier than currently priced by markets.

\begin{table}[htbp]
\caption{EUR bond ETF universe}
\footnotesize
\begin{tabularx}{\linewidth}{P{1.75cm} P{1.0cm} P{2.35cm} P{1.5cm} Q{0.8cm} Q{1.15cm} Q{0.95cm} L}
\toprule
\tabh{Bucket} & \tabh{Ticker} & \tabh{ISIN} & \tabh{Category} & \tabh{TER} & \tabh{Duration (yrs)} & \tabh{Weight (\%)} & \tabh{Function} \\
\midrule
Short (50\%) & XEON & LU0290358497 & Money Market & 0.10\% & \(\sim\)0.0 & 15\% & Tracks overnight ECB rate ({\texteuro}STR), minimal price risk \\
 & ERNE & IE00BCRY6557 & Ultrashort / Floater & 0.09\% & \(\sim\)0.4 & 8\% & Mix of very short fixed-rate and floating-rate IG bonds \\
 & EFRN & IE00BF5GB717 & Floater & 0.10\% & \(\sim\)0.2 & 7\% & Pure floating-rate corporate sleeve, resets on Euribor \\
 & DBXP & LU0290356871 & Govt 1--3yr & 0.10\% & \(\sim\)1.9 & 8\% & Core short government bond position, iBoxx index \\
 & SXRN & IE00B3VTMJ91 & Govt 1--3yr & 0.15\% & \(\sim\)1.9 & 7\% & Provider diversification via Bloomberg index \\
 & XYP1 & LU0925589839 & Govt 1--3yr & 0.15\% & \(\sim\)1.9 & 5\% & Tilt toward higher-yielding Eurozone sovereigns for carry \\
\midrule
Intermediate (15\%) & EUNT & IE00B4L60045 & IG Corporate & 0.20\% & \(\sim\)2.8 & 5\% & Investment grade corporate spread exposure, distributing \\
 & ICOV & IE00B3B8Q275 & Covered Bonds & 0.20\% & \(\sim\)4.5 & 4\% & Secured bank bonds with very high credit quality \\
 & IEAC & IE00B3F81R35 & IG Corporate & 0.09\% & \(\sim\)4.7 & 6\% & Broad IG corporate core across all maturities, lowest TER \\
\midrule
Long (35\%) & IBGX & IE00B1FZS681 & Govt 3--5yr & 0.15\% & \(\sim\)3.9 & 8\% & Transition step toward longer duration \\
 & SXRP & IE00B3VTML14 & Govt 3--7yr & 0.15\% & \(\sim\)4.8 & 9\% & Extension of government duration anchor \\
 & IBCL & IE00B1FZS913 & Govt 15--30yr & 0.15\% & \(\sim\)15.9 & 18\% & Primary long-duration barbell anchor; equity diversification and growth shock protection \\
\bottomrule
\end{tabularx}
\end{table}

The long-duration allocation within the EUR sleeve is anchored primarily by IBCL. While this product carries meaningful duration risk in isolation, it serves a clearly defined role within the barbell framework: providing the portfolio with convexity, flight-to-quality upside in a growth shock scenario and the diversification benefit against equity markets that Section 4.1 identifies as a core portfolio objective. IBCL alone contributes roughly 62 percent of the EUR sleeve's duration: on the table's weights and durations, the EUR sleeve runs approximately 4.6 years, the USD sleeve approximately 3.5 years, and the blended 70/30 book approximately 4.3 years, which is the figure behind the portfolio's lower-beta positioning relative to long-duration benchmarks.

\subsection{USD Products (30\% of total portfolio)}

The USD sleeve follows a differently weighted barbell structure --- 20\% short duration, 50\% intermediate duration and 30\% long duration, as set out in Section 3.2 --- adapted to the specific dynamics of the US rate environment described in Section 2.3. The critical distinction from the EUR sleeve is the role of the 1--3 year segment. While short-dated Eurozone government bonds represent a relatively stable short-duration position, the equivalent US Treasury segment absorbed the sharpest yield repricing following the June 2026 FOMC meeting, as markets adjusted to a higher projected terminal rate. IUSU is therefore included but held at a deliberately reduced weight within the short bucket, with money market and floating rate instruments taking a proportionally larger share. The long-duration component benefits from a structural anchor that is particularly strong on the USD side: the Fed's longer-run neutral rate projection of approximately 3.0\% provides a fundamental support for the long end of the Treasury curve, making the long-duration barbell position a well-supported structural exposure consistent with the Goldman Sachs assessment referenced in Section 2.3.2.3, which projects rate cuts materializing from mid to end 2027. Two implementation notes apply to this sleeve. The XYLD line should be entered by ISIN (IE00BF8J5974) rather than by ticker, because the bare ticker XYLD also designates a US-listed equity covered-call fund and ordering by ticker risks buying the wrong instrument, and the fund's modest size of roughly EUR 71 million warrants a liquidity check at real sizing.

\begin{table}[htbp]
\caption{USD bond ETF universe}
\footnotesize
\begin{tabularx}{\linewidth}{P{1.75cm} P{1.0cm} P{2.35cm} P{1.5cm} Q{0.8cm} Q{1.15cm} Q{0.95cm} L}
\toprule
\tabh{Bucket} & \tabh{Ticker} & \tabh{ISIN} & \tabh{Category} & \tabh{TER} & \tabh{Duration (yrs)} & \tabh{Weight (\%)} & \tabh{Function} \\
\midrule
Short (20\%) & IB01 & IE00BGSF1X88 & Treasury 0--1yr & 0.07\% & \(\sim\)0.4 & 6\% & Near-cash T-Bill proxy, minimal price risk \\
 & ERND & IE00BCRY6227 & Ultrashort IG & 0.09\% & \(\sim\)0.4 & 6\% & USD equivalent of ERNE \\
 & FLOA & IE00BDFGJ627 & Floater & 0.10\% & \(\sim\)0.0 & 5\% & Floating rate corporates resetting on SOFR \\
 & IUSU & IE00B14X4S71 & Treasury 1--3yr & 0.07\% & \(\sim\)1.9 & 3\% & Held at reduced weight given front-end repricing risk \\
\midrule
Intermediate (50\%) & SDIG & IE00BCRY5Y77 & IG Corporate & 0.20\% & \(\sim\)2.6 & 32\% & USD equivalent of EUNT, short-duration IG spread exposure \\
 & XYLD & IE00BF8J5974 & IG Corporate (ESG) & 0.16\% & \(\sim\)2.4 & 18\% & Second-provider USD IG credit exposure, diversifying the intermediate bucket beyond a single index/provider \\
\midrule
Long (30\%) & CBU7 & IE00B3VWN393 & Treasury 3--7yr & 0.07\% & \(\sim\)4.6 & 12\% & Step into longer duration, stable relative to short end \\
 & CBU0 & IE00B3VWN518 & Treasury 7--10yr+ & 0.07\% & \(\sim\)8.5 & 18\% & Primary long barbell anchor; benefits from longer-run Fed rate normalisation from 2027 \\
\bottomrule
\end{tabularx}
\end{table}

\section{Conclusion and Monitoring}

\subsection{Summary}

The Eurozone and US analyses point to two distinct monetary policy environments: the ECB is tightening into weak growth as inflation risks remain elevated, while the Fed is expected to hold rates broadly stable through 2026 before easing gradually from 2027 onwards. This divergence underpins the different sleeve structures adopted in Section 4, a mild barbell for the EUR sleeve and a bullet-weighted, intermediate-focused structure for the USD sleeve, and is reflected in the 70\%/30\% currency split and the 50/15/35 and 20/50/30 duration allocations respectively. Together, these choices satisfy the three guiding principles set out in Section 1.3: a defensive, shock-absorbing role, a lower-beta profile relative to long-duration benchmarks, and diversification across currencies and maturity segments. The portfolio deliberately avoids large directional bets on the ECB's or the Fed's next move, prioritising capital preservation and stability over yield maximisation.

\subsection{Limitations}

Two limitations should be kept in mind when interpreting this report. First, as noted in Section 1.1, the analysis excludes political and geopolitical developments after 26 June 2026, including the renewed escalation of the Iran conflict; any material shift in energy prices or risk sentiment since that date is not reflected in the positioning above. Second, market-implied pricing referenced throughout this report, such as the roughly one further ECB hike priced for 2026 or the Fed cuts expected from mid-2027, reflects a snapshot of market expectations rather than a forecast, and can shift materially as new data arrives.

\subsection{Monitoring Framework}

Ongoing monitoring should focus on five areas. On the ECB side, this includes Governing Council meeting outcomes, the evolution of core versus headline inflation, negotiated wage growth and energy price developments. On the Fed side, the key inputs are FOMC meetings and Summary of Economic Projections updates, PCE and Sticky Price CPI, and labour market indicators such as unemployment, payrolls and wage growth. Eurozone sovereign spreads, particularly Italy and France versus Bunds, should be tracked alongside any change in country debt ratios. The EUR--USD interest rate differential should be monitored as the primary driver of hedging costs for the USD sleeve. Finally, developments in the Middle East energy shock, and the Russia--Ukraine conflict in the background, remain the key geopolitical variables for both currency areas.

\subsection{Rebalancing Triggers}

Specific triggers for rebalancing, as anticipated in Section 1.3, are as follows. In the EUR sleeve, ECB hikes beyond current market pricing should prompt a further shortening of duration, while clear signs of recession should shift weight toward the long bucket. In the USD sleeve, Fed easing signals materialising earlier than mid-2027 should increase the intermediate and long weights, while more persistent inflation should prompt further shortening. On the currency side, a material narrowing of the EUR--USD rate differential should prompt a reassessment of the currently unhedged USD position. Sovereign spread widening beyond a defined threshold should trigger a reduction in peripheral exposure within the EUR short bucket. Independently of these thresholds, the portfolio should undergo a fixed quarterly rebalancing review to confirm that sleeve and bucket weights remain within their target ranges.

Taken together, these monitoring points and triggers are intended to keep the portfolio anchored to its original mandate: a safe, low-beta foundation within a broader multi-asset context, even as the interest rate environment in both currency areas continues to evolve.

\clearpage

\begin{fshfreferences}
\addcontentsline{toc}{section}{References}
% Allow URL line breaks at hyphens within this list only (scoped to the environment group)
\expandafter\def\expandafter\UrlBreaks\expandafter{\UrlBreaks\do\-}
\item CME Group (n.d.). CME FedWatch Tool. \url{https://www.cmegroup.com/markets/interest-rates/cme-fedwatch-tool.html}
\item ECB Watch (n.d.). ECB Interest Rate Expectations and Monetary Policy Monitoring. \url{https://ecb-watch.eu/}
\item European Central Bank (n.d.). ECB Data Portal. \url{https://data.ecb.europa.eu/}
\item European Central Bank (n.d.). ECB Wage Tracker (EWT) Dataset. \url{https://data.ecb.europa.eu/data/datasets/EWT}
\item European Central Bank (2026). Monetary Policy Decisions, 11 June 2026. \url{https://www.ecb.europa.eu/press/pr/date/2026/html/ecb.mp260611~4d41bd5e83.en.html}
\item European Commission, Eurostat (n.d.). Euro Indicators. \url{https://ec.europa.eu/eurostat/news/euro-indicators}
\item European Commission, Eurostat (2026). GDP down by 0.2\% and employment up by 0.1\% in the euro area. Euro Indicators. \url{https://ec.europa.eu/eurostat/web/products-euro-indicators/w/2-05062026-ap}
\item Federal Reserve Bank of New York (n.d.). Effective Federal Funds Rate (EFFR). \url{https://www.newyorkfed.org/markets/reference-rates/effr}
\item Federal Reserve Bank of St.\ Louis (n.d.). FRED --- Federal Reserve Economic Data. \url{https://fred.stlouisfed.org}
\item Federal Reserve Bank of St.\ Louis (n.d.). FRED graph: core labor market data. \url{https://fred.stlouisfed.org/graph/?g=1WWDy}
\item Federal Reserve Bank of St.\ Louis (n.d.). FRED series PCEPILFE: core inflationary data, percent change from year ago. \url{https://fred.stlouisfed.org}
\item Federal Reserve Board (n.d.). Monetary Policy. \url{https://www.federalreserve.gov/aboutthefed/fedexplained/monetary-policy.htm}
\item Federal Reserve Board (n.d.). PCE Inflation. \url{https://www.federalreserve.gov/economy-at-a-glance-inflation-pce.htm}
\item Federal Reserve Board, Federal Open Market Committee (2026). Summary of Economic Projections (FOMC ``Dot Plot''), 17 June 2026.
\item Federal Statistical Office of Germany (Destatis) (n.d.). Economy and Finance --- Key Indicators for Europe. \url{https://www.destatis.de/Europa/EN/Topic/Key-indicators/EconomyFinance.html}
\item Goldman Sachs (n.d.). Research report revising US interest rate cut expectations from end-2026 to mid- and end-2027 (as cited in text; no publication details given).
\item National Bank of Belgium (2026). Monetary Policy Decisions, June 2026. \url{https://www.nbb.be/en/news-events/news/monetary-policy-decisions-june-2026}
\item Rate Probability (n.d.). ECB Rate Probability Tracker. \url{https://rateprobability.com/ecb}
\end{fshfreferences}

\end{document}
